% concatenated from main_2.tex
%\documentclass{article}
\documentclass[12pt]{amsart}
\usepackage{amsfonts, amsbsy, amsmath, amssymb}
\usepackage{xcolor}
\usepackage{graphicx} % Required for inserting images
\newcommand\Copi{\rotatebox[origin=c]{180}{$\Pi$}}
\usepackage{amsmath}
\usepackage{amssymb}
\usepackage{tikz-cd}


\usepackage{physics}
\newcommand{\Leib}{\mathrm{Leib}}




\theoremstyle{plain}
\newtheorem*{theorem*}{Theorem}
\newtheorem*{thmex*}{Theorem~\ref{example}}
\newtheorem*{thmasymp*}{Theorem~\ref{thmAsymp}}
\newtheorem{theorem}{Theorem}[section]
\newtheorem{acknowledgement}[theorem]{Acknowledgement}
\newtheorem{algorithm}[theorem]{Algorithm}
\newtheorem{assertion}[theorem]{Assertion}
\newtheorem{axiom}[theorem]{Axiom}
\newtheorem{case}[theorem]{Case}
\newtheorem{claim}[theorem]{Claim}
\newtheorem{conclusion}[theorem]{Conclusion}
\newtheorem{condition}[theorem]{Condition}
\newtheorem{conjecture}[theorem]{Conjecture}
\newtheorem{corollary}[theorem]{Corollary}
\newtheorem{criterion}[theorem]{Criterion}
\newtheorem{exercise}[theorem]{Exercise}
\newtheorem{lemma}[theorem]{Lemma}
\newtheorem{property}[theorem]{Property}
\newtheorem{notation}[theorem]{Notation}
\newtheorem{problem}[theorem]{Problem}
\newtheorem{proposition}[theorem]{Proposition}
\newtheorem{remark}[theorem]{Remark}
\newtheorem{scholium}[theorem]{Scholium}
\newtheorem{solution}[theorem]{Solution}
\newtheorem{summary}[theorem]{Summary}
\newtheorem{example}[theorem]{Example}

\usepackage{comment}

\usepackage{datetime}

\theoremstyle{definition}
\newtheorem{definition}[theorem]{Definition}


\usepackage[lite]{amsrefs}
\newcommand*{\MRref}[2]{ \href{http://www.ams.org/mathscinet-getitem?mr=#1}{MR \textbf{#1}}}
\newcommand*{\arxiv}[1]{\href{http://www.arxiv.org/abs/#1}{arXiv: #1}}
\renewcommand{\PrintDOI}[1]{\href{http://dx.doi.org/\detokenize{#1}}{doi: \detokenize{#1}}%
	\IfEmptyBibField{pages}{, (to appear in print)}{}}


\title{Lie Quandles, Leibniz Racks and Noether's First Theorem}

\author{Mohamed Elhamdadi}
\address{Department of Mathematics and Statistics, University of South Florida, Tampa, Florida, USA}
\email{emohamed@usf.edu}

\author{Bryce Virgin}
\address{Department of Mathematics and Statistics, University of South Florida, Tampa, Florida, USA}
\email{bgvirgin@usf.edu}
\date{April 17, 2026}
%\today


\begin{document}

\begin{abstract}

%Quandles are algebraic constructs that have arisen from both the study of knots, and from the study of the conjugation structure of groups. They are a concrete realization of the desire to form an algebraic object whose elements determine automorphisms of itself Recently, T. Fritz in his paper "Self-distributive structures in physics" indicated how quandles arise in the context of the Hamiltonian formulation of classical and quantum mechanics, as the evolution structure on the algebra of observables determined by the Lie bracket and a choice of Hamiltonian observable. In particular, Fritz indicated that his observation leads to a natural nonlinear generalization which he dubbed "Lie Quandles", which realize the desire to form an algebraic object whose elements determine smooth one parameter families of autmorphisms of itself. 
In [Self-distributive structures in physics.
Internat. J. Theoret. Phys. 64 (2025), no. 3, Paper No. 73], Fritz was motivated by the structure of Hamiltonian/Heisenberg mechanics to define the notion of "Lie Quandle", which he argued are nonlinear generalizations of finite dimensional real Lie algebras. In this article, we will investigate a  linear/nonlinear correspondence to which Fritz' is a special case, classify a 
class of generalizations of these objects, as well as describe some results in the direction of a nonlinear analogue of Noether's first theorem first described by Fritz.

\end{abstract}

\maketitle

\tableofcontents









\section{Introduction}\label{intro}
%%%%%%%%%%%%%%%
Quandles and Racks are algebraic constructs that have arisen from both the study of knots \cite{Joyce, Matveev, EN}, and from the study of the conjugation structure of groups \cite{BES}. Racks are a concrete realization of the desire to form an algebraic object whose elements can be assigned to automorphisms of itself %(check if this is still the correct reference) 
\cite{Brieskorn}. Quandles appeared in the literature with many different names.  Mituhisa Takasaki \cite{Takasaki} introduced, in 1942, the notion of kei as an abstraction of the notion of symmetric transformation.
The earliest known work on racks is contained in the 1959 correspondence (unpublished) between John Horton Conway and Gavin Wraith who studied
racks in the context of the conjugation operation in a group.  Around 1980, Joyce \cite{Joyce} and Matveev \cite{Matveev} independently introduced the notion of a quandle. To each oriented knot, they associated a quandle in such a way that the classical problem of determining knot equivalence is translated into the problem of quandle isomorphism. 

Since their introduction, quandles have been the subject of extensive study in both topology and algebra. From an algebraic perspective, they are naturally related to a range of algebraic structures, including Lie algebras \cite{CCES1, CCES2},  Frobenius algebras \cite{CCEKS} and the Yang-Baxter equation \cite{BES}, Hopf algebras \cite{And-Grana, CCES2}, quasigroups and Moufang loops \cite{Elhamdadi}, and ring theory \cite{EFT}. 
Quandles have also been used in the construction of invariants for handlebody knots by enriching their structure with a group-indexed family of operations, leading to the notion of a $G$-family of quandles introduced in \cite{IIJO}. This enhanced framework allows for the development of stronger invariants for knots, links, and related objects such as spatial graphs. 

Topological quandles were investigated in \cite{Rubinsztein} and their cohomology was studied in \cite{EM, ESZ}.  Quandles and Lie racks, objects partially solving the coquecigrue problem, were investigated in \cite{Kinyon}. In \cite{Covez} the author showed that
Leibniz algebras integrate (locally) to local Lie racks, providing a precise solution to the coquecigrue problem \cite{Loday}. In \cite{LucTa} the author showed that affine algebraic racks provide a solution to a version of the well-known coquecigrue problem for left and right Leibniz algebras over arbitrary fields.  Taking a family of quandles at a time, the notion of $G$-family of quandles was introduced in \cite{IIJO} with the goal of producing invariants of handlebody knots.   

%{\color{red} (Discuss Fritz' paper, his motivations for defining Lie quandles, and his conjectured form of Noether's first theorem.) (Discuss the fact that the space of observable really ought to be allowed to be infinite dimensional)}

Tobias Fritz introduced the notion of Lie quandle in \cite{Fritz} as a nonlinear generalization of finite dimensional real Lie algebras. Fritz was motivated to consider this generalization by the way observables in both the Hamiltonian formulation of classical mechanics and the Heisenberg formulation of quantum mechanics generate 1-parameter subgroups of automorphisms of the entire observable algebra. This feature of these frameworks led to a particularly simple way of stating Noether's first theorem, which necessarily leads to an equivalent statement of Noether's first theorem for Lie quandles. The Lie Quandles section in \cite{Fritz} concludes with some examples of Lie quandles both satisfying Noether's first theorem and not, and a suggestion that connectedness may be a sufficient condition to imply Noether's first theorem on a Lie quandle. 



In this article, we describe some of the relationships between some categories related to Lie quandles and Lie algebras, discuss a special case classification result of one of these categories, and investigate the necessity or sufficiency of possible hypotheses for a further generalizations of Noether's first theorem in addition to Fritz' suggestion of connectedness.

The paper is organized as follows:  In Section~\ref{Sec2}, we review basic definitions, including those of Lie quandles and 
$G$-families of quandles, and introduce new terminology, such as the notion of a smooth $G$-family of smooth racks. In Section~\ref{Sec3} we show that finite dimensional real Leibniz algebras have the same relationship to smooth $\mathbb{R}$-families of smooth racks that finite dimensional real vector spaces have to smooth manifolds, and consequently this correspondence is true of finite dimensional real Lie algebras and Lie quandles as suggested by Fritz' argument that the Lie quandle naturally associated to a Lie algebra by exponentiation contains all of the information of the originial Lie algebra. We also show that smooth $G$-families of smooth racks ($G$ a simply connected Lie group) can be reduced to certain families of smooth $\mathbb{R}$-families of smooth racks. Section~\ref{Sec4} provides specialized classification of a simple class of smooth $G$-families of smooth racks. We conclude with Section~\ref{Sec5}, in which we discuss Noether's first theorem in the context of smooth $G$-families of smooth racks, and among these results we see that connectedness is not a necessary condition.




\vspace{5mm}

\section{Review of racks and Leibniz algebras}\label{Sec2}
We start by briefly recalling some definitions and examples.
%%%%%%%%%%%%%
\begin{definition}
     A rack is a pair $(Q,\triangleleft)$ where $Q$ is a set and $\triangleleft:Q^2 \to Q$ is a map satisfying the following two axioms:

\begin{enumerate}

    \item 
 The right multiplication $R_x: Q \to Q, y \mapsto y \triangleleft x$, is a bijection,

\item 
For all $p,q,r \in Q$, $ (p  \triangleleft q )\triangleleft r = (p  \triangleleft r )\triangleleft (q  \triangleleft r).$ 
\end{enumerate}

\end{definition}
If $p  \triangleleft p=p$ holds for all $p \in Q$, then $(Q,\triangleleft)$ is called a quandle.\\

 Typical examples of quandles include the following. 
  \begin{itemize}

  \item
Any non-empty set $Q$ with the operation $p \triangleleft q=p$ for any $p,q \in Q$ is
a quandle called a  {\it trivial} quandle.

\item
A conjugacy class $Q$ of a group $G$ with the quandle operation $p \triangleleft q=q^{-1}pq $ is called a {\it conjugation quandle}. 

\item
For a group $G$, the binary operation $p \triangleleft q=q p^{-1}q $ is called a {\it core quandle}. 


  \item 
  A {\it generalized Alexander quandle} is defined  by 
a pair $(G, f)$ where 
$G$ is a  group and $f \in {\rm Aut}(G)$,
and the rack operation is defined by 
$p \triangleleft q=f(pq^{-1}) q $. 
If $G$ is abelian, this is called an {\it Alexander quandle}. 

  

\end{itemize}
%%%%%%%%%%%%%

\begin{definition}
     A Lie quandle is a pair $(Q,\triangleleft)$ where $Q$ is a smooth manifold and $\triangleleft:Q^2\times \mathbb{R} \to Q$ is a smooth map having the following three properties:

\begin{enumerate}
    \item 
 For all $q \in Q$ and $ t \in \mathbb{R}$, $q \triangleleft_t q =q,$  

\item 
For all $q,p \in Q$ and $t,s \in \mathbb{R}$, $(q  \triangleleft_t p )\triangleleft_s p = q \triangleleft_{t+s} p $ and $q\triangleleft_0 p = q$,

\item 
For all $ t,s \in \mathbb{R}$  
and $q,p,r \in Q$, $(q\triangleleft_t p) \triangleleft_s r = (q \triangleleft_s r)\triangleleft_{t} ( p\triangleleft_s r) $ 
\end{enumerate}

\end{definition}


\vspace{5mm}

This definition turns out to be a specialization of an object, called a $G$-family of quandles \cite{IIJO}, which has already found its way into the quandle literature, corresponding to collections of quandles parameterized by group elements.

\vspace{5mm}

\begin{definition}
    

A $G$-family of quandles for a given group $G$ is a pair $(Q,\triangleleft)$ where $Q$ is a set and $\triangleleft:Q^2\times G \to Q$ is a function satisfying:
\begin{enumerate}
    \item 
 $q \triangleleft_g q =q$ for all $q \in Q$ and $g \in G$.

\item 
$ (q  \triangleleft_g p )\triangleleft_h p = q \triangleleft_{gh} p $ and $q\triangleleft_e p = q$ for all $q,p \in Q$.

\item 
$ (q\triangleleft_g p) \triangleleft_h r = (q \triangleleft_h r)\triangleleft_{h^{-1}gh} (p\triangleleft_h r) $ for all $g,h \in G$ and $q,p,r \in Q$.
\end{enumerate}
\end{definition}

We will denote the $G$-family of quandles by $(Q, \{\triangleleft_g\}_{g \in G})$.  To make the notation shorter we will use $(Q,G)$ instead. 
\vspace{5mm}

The following are a few examples of $G$-families of quandles.
    \begin{itemize}
        \item 
        For any group $G$ and any set $Q$, defining $p \triangleleft_g q=p$ for all $p,q \in Q$ and all $g \in G$.  This gives a $G$-family of quandles called the \emph{trivial} $G$-family of quandles.
        
        \item
        Let $(Q, \triangleleft )$ be a quandle of cardinality $n+1$, such that ${R_q}^{n}(p)=p$ for all $p,q \in Q$.  Define $p \triangleleft_iq={R_q}^i(p)$ then  $(Q,\mathbb{Z})$ is a $\mathbb{Z}$-family of quandles and also $\mathbb{Z}_{n}$-family of quandles.

        \item 
        For a ring $R$ and a group $G$ let $R[G]$ be the group ring.  Then any (right) $R[G]$-module $Q$ becomes a $G$-family of quandles by defining $p \triangleleft_g q:=pg+q(1-g)$.
    \end{itemize}

    Notice that a $G$-family of quandles $(Q,G)$ induces a quandle operation on $Q \times G$ given by 
    \[
    (p,g)*(q,h):= (p \triangleleft_g q, h^{-1}gh). 
    \]

 We are grateful to Masahico Saito for observing that Lie quandles can be viewed as 
instances of $\mathbb{R}$-families of quandles internal to the category of smooth manifolds. Given this observation, it is natural to internalize the general notion of $G$-families of quandles to the category of smooth manifolds, and to subject these new objects to Fritz' line of questioning.

\vspace{5mm}

\begin{definition}
     We define a \textbf{smooth} $G$-\textbf{family of smooth racks} for a given (real) Lie group $G$ to be a pair $(R,\triangleleft)$ where $R$ is a smooth manifold and $\triangleleft:R^2\times G \to R$ is a smooth map having the following three properties:

\begin{enumerate}
\item 
$q\triangleleft_e p = q$ for all $q,p \in R$.
    \item 
 $ (q  \triangleleft_g p )\triangleleft_h p = q \triangleleft_{gh} p$ for all $g,h \in G$ and $q,p \in R$.

\item 
$ (q\triangleleft_g p) \triangleleft_h r = (q \triangleleft_h r)\triangleleft_{h^{-1}gh} (p\triangleleft_h r) $ for all $g,h \in G$ and $q,p,r \in R$.
\end{enumerate}
\end{definition}
A smooth $G$-family of smooth racks $(R,\triangleleft)$ is called a smooth $G$-family of smooth quandles if $q\triangleleft_g q =q$ for all $q \in R$ and $g \in G$. A smooth $\mathbb{R}$-family of smooth racks is called a \textbf{Leibniz Rack}. 

\begin{example}
    For any smooth manifold $R$ and any Lie group $G$, we can define the trivial smooth $G$-family of smooth rack operations via $p \triangleleft_g q=p$ for all $p,q \in R$ and all $g \in G$. If $G  =\mathbb{R}$, we call this family of operations the trivial Leibniz rack operation for $R$.
\end{example}

\begin{example}
    All Lie quandles are Leibniz racks.
\end{example}


\begin{example}
    For any Lie groups $F,G$, any closed Lie subgroup $H \subset F$, and any Lie group action $g \mapsto \rho_g \in \text{Aut}(F)$ such that $\rho_g(H) \subset H$ for all $g \in G$, the operation $f'H \triangleleft_g fH = f\rho_g(f^{-1}f')H$ is a smooth $G$-family of smooth rack operations on the coset manifold $F/H$. In ~\ref{Sec5} we will see that a special subset of this family of smooth $G$-families of smooth racks obeys Noether's first theorem.
\end{example}


\vspace{5mm}

\begin{definition}
    A morphism between two smooth Lie group families of smooth racks $(R,\triangleleft)$ and $(R',\triangleleft')$ is a smooth map $f:R\to R'$ and a Lie group homomorphism $\varphi:G \to G'$ such that:
    $$ f(p \triangleleft_g q) = f(p) \triangleleft_{\varphi(g)}'f(q) $$
\end{definition}

\vspace{5mm}


Due to the correspondence between simply connected Lie groups and finite-dimensional real Lie algebras, the category of smooth $G$-families of smooth racks for a simply connected $G$ is isomorphic to the category whose objects are families of Leibniz racks parameterized properly by the elements of a Lie algebra, as we will see in section ~\ref{Sec3}. 

 \vspace{5mm}

 Leibniz algebras were introduced by Loday \cite{Loday} as a non-commutative version of Lie algebras, where the bracket is not necessarily antisymmetric.  We recall the definition.

%\textbf{Definition}
\begin{definition}    
A (right) Leibniz algebra is a vector space $A$ together with a bilinear bracket $[\cdot,\cdot]$ which has the (right) Leibniz property:
$$ [[a,b],c] = [[a,c],b]+[a,[b,c]]. $$
Every Lie algebra is a Leibniz algebra, and for both Lie and Leibniz algebras we adopt the convention that for each $a \in A$ the map $\text{ad}_a:A \to A$ is the inner derivation defined $\text{ad}_a(b) = [b,a]$ for each $b \in A$.
\end{definition}

\begin{example}
  It's clear that a Lie algebra is an example of Leibniz algebra.  In particular, $\mathfrak{sl}_2(\mathbb{R})$ with bracket $[X,Y]=XY-YX$ is a Leibniz algebra.  
\end{example}

\begin{example}[A non-Lie Leibniz algebra]
Let $L = \mathrm{span}\{e_1,e_2\}$ with
$[e_1,e_1] = e_2,$ and all other brackets being zero.
Then $L$ is a Leibniz algebra but not a Lie algebra.
\end{example}

\begin{example}[From associative algebras]
Let $A$ be an associative algebra. Define
\[
[x,y] := xy-yx.
\]
Then the Leibniz identity follows from associativity.
\end{example}

\begin{example}[Hemisemidirect product]
Let $\mathfrak{g}$ be a Lie algebra acting on a vector space $V$. Define
\[
[(x,v),(y,w)] = ([x,y], x \cdot w - y \cdot v).
\]
This gives a Leibniz algebra structure on $\mathfrak{g} \oplus V$.
\end{example}



%\subsection{The Ideal of Squares}

A fundamental structural feature of Leibniz algebras is the following.

\begin{proposition}
Let $L$ be a Leibniz algebra and define
\[
\Leib(L) = \langle [x,x] : x \in L \rangle.
\]
Then $\Leib(L)$ is a two-sided ideal of $L$, and the quotient $L/\Leib(L)$ is a Lie algebra.
\end{proposition}

\begin{example}
For a Leibniz algebra $L$ defined by $[e_1,e_1]=e_2$, we have
\[
\Leib(L) = \mathrm{span}\{e_2\},
\]
and $L/\Leib(L)$ is abelian.
\end{example}

%\subsection{Low-Dimensional Classification}

%Low-dimensional Leibniz algebras have been classified and provide useful intuition.

%\subsubsection{Two-Dimensional Case}

Let $L$ be a $2$-dimensional Leibniz algebra over $\mathbb{K}$. Then, up to isomorphism, the non-Lie case is given by:
\[
[e_1,e_1] = e_2,\; \text{and} \;\; [e_1,e_2]=[e_2,e_1]=[e_2,e_2]=0
\]

All other $2$-dimensional Leibniz algebras are Lie algebras.

%%%%%%%%%
\begin{comment}
    
%\subsubsection{Three-Dimensional Examples}
\begin{example}

In dimension three, Leibniz algebras are no longer unique up to isomorphism; indeed, there exist several non-isomorphic structures \cite{AO}. Typical examples include:

\begin{itemize}
\item $[e_1,e_1] = e_2$
\item $[e_1,e_2] = e_3$
\item $[e_1,e_1] = e_3,\quad [e_2,e_2] = e_3$
\item $[e_1,e_2] = e_2,\quad [e_2,e_1] = -e_2$
\end{itemize}
\end{example}
These illustrate the richness of the non-antisymmetric structure.
\end{comment}





%%%%%%%%%%%%%%
%%%%%%%%%%%%%%
\vspace{5mm}

%\textbf{Definition}
\begin{definition}
    We define a smooth $A$-family of Leibniz racks for a given finite dimensional real Leibniz algebra $A$ to be a pair $(R,\triangleleft )$ where $R$ is a smooth manifold and $\triangleleft:R^2\times A\times \mathbb{R} \to R$ is a smooth map (with evaluation denoted $(q,p,a,t) \mapsto q\triangleleft_{(a,t)}p$, and we also define the operation $\triangleleft^a:R^2 \times \mathbb{R} \to R$ via $q\triangleleft^a_t p = q\triangleleft_{(a,t)}p$ for all $t \in \mathbb{R}$, $a \in A$, and all $q,p \in R$) satisfying the following two axioms:

\begin{enumerate}
    \item 
    For every $a \in A$, the pair $(R,\triangleleft^a)$ is a Leibniz rack.

    \item For every $a \in A$, every $t \in \mathbb{R}$, and all $p,q \in R$, it follows that $ p \triangleleft_t^a  q =p \triangleleft_{1}^{ta} q$.

\item  For all $a,b \in A$, all $t,s \in \mathbb{R}$, and all $o,p,q \in R$, we have that:
$$ (q \triangleleft_t^a p) \triangleleft_s^b o = (q  \triangleleft_{s}^b o ) \triangleleft_{ t }^{e^{s \cdot \text{ad}_b}(a)} (p  \triangleleft_{s}^b o ). $$

\end{enumerate}
\end{definition}



\begin{definition}
  A morphism between two smooth Leibniz families of Leibniz racks $(R, \triangleleft   )$ and $(R',\triangleleft')$ consists of a smooth map $f:R\to R'$ and a Leibniz algebra morphism $\varphi:A \to A'$ such that: $$f( p \triangleleft_{(a,t)} q) = f(p) \triangleleft'_{(\varphi(a),t)} f(q) \text{ for all }q,p \in R, \text{ }a\in A, \text{ and }t \in \mathbb{R}.$$
\end{definition}


\section{Leibniz Algebras and Leibniz Racks} \label{Sec3}



In \cite{Fritz}, Fritz showed how the Lie bracket of a Lie algebra may be recovered from the natural Lie quandle operation associated to it. His explanation formed the core of an argument that Lie quandles have the same relationship to finite dimensional real Lie algebras that smooth manifolds have to finite dimensional real vector spaces. This turns out to be a special case of a more general correspondence between finite dimensional real Leibniz algebras and Leibniz racks. The proof of the following lemma likely exists elsewhere in the literature, we simply include it for the benefit of this article's self-containment.

\vspace{5mm}

 

\begin{lemma}
     For any finite-dimensional, real Leibniz algebra $A$, and any $a \in A$, we have that $e^{ \text{ad}_a} :A \to A $ is an automorphism of Leibniz algebras.
     \end{lemma}
     
\begin{proof}
 We begin by defining the smooth functions $f,g : \mathbb{R}\to A$ via:
$$ f(t) = e^{t [\cdot,a]}([b,c]) $$
$$ g(t) = [e^{t [\cdot,a]}(b),e^{t [\cdot,a]}(c)] $$
We note that $f(0) = [b,c] = g(0)$, and furthermore that:
$$ f'(t) = [e^{t [\cdot,a]}([b,c]),a] = [f(t),a] $$
$$ g'(t) = [[e^{t [\cdot,a]}(b),a],e^{t [\cdot,a]}(c)]+[e^{t [\cdot,a]}(b),[e^{t [\cdot,a]}(c),a]] $$
$$= [[e^{t [\cdot,a]}(b),e^{t [\cdot,a]}(c)],a]  = [g(t),a] $$
hence $f$ and $g$ are both solutions to the same first order linear initial value problem (IVP) $u' = [u,a]$ with $u(0) = [b,c]$, hence by the uniqueness of solutions to first order linear IVPs we must have that $f=g$. By setting $t=1$, we learn that: $$ e^{ [\cdot,a]}([b,c]) = [e^{t [\cdot,a]}(b),e^{t [\cdot,a]}(c)],$$
furthermore, $e^{[\cdot,a]}$ has $e^{-[\cdot,a]}$ as an inverse, and therefore $e^{[\cdot,a]}$ is a Leibniz algebra automorphism.
    
\end{proof}

As a simple corollary to the lemma above, we have that for any Leibniz algebra $A$ and any $a,b,c \in A$ we have that:
$$ (e^{\text{ad}_a} \circ \text{ad}_b \circ e^{-\text{ad}_a})(c) = e^{\text{ad}_a}([e^{-\text{ad}_a}(c),b]) = [c,e^{\text{ad}_a}(b)] = \text{ad}_{e^{\text{ad}_a}(b)}(c) $$
$$ \implies e^{\text{ad}_a} \circ \text{ad}_b \circ e^{-\text{ad}_a} = \text{ad}_{e^{\text{ad}_a}(b)}, $$
a fact which will turn out to be useful for the next proposition.


\begin{proposition}
The category of Leibniz algebras is isomorphic to a subcategory of the category of Leibniz racks whose morphisms are precisely all linear maps. The functor constructed to show this restricts to Fritz's isomorphism of the category of (real, finite dimensional) Lie algebras to the subcategory of Lie quandles whose morphisms are precisely all linear maps.
\end{proposition}

\begin{proof}
 For any finite dimensional real Leibniz algebra $A$, define $F(A)$ to the pair $(A,\triangleleft)$ where $a\triangleleft_s b = e^{s [\cdot,b]}(a)$ for all $a,b \in A$ and $s \in \mathbb{R}$. Since we'll need it later, we note that: $$ \frac{d}{ds}\bigg|_{s=0} (a\triangleleft_s b) = \frac{d}{ds}\bigg|_{s=0} e^{s [\cdot,b]}(a) = [a,b].$$

Now to see that $F(A)$ is a smooth $\mathbb{R}$-family of smooth racks, we first check that:
$$  a\triangleleft_0 b = e^{0 [\cdot,b]}(a) = a  $$
while at the same time:
$$ (a \triangleleft_t b ) \triangleleft_s b  = e^{s [\cdot,b]}e^{t [\cdot,b]}a = e^{(t+s) [\cdot,b]}a = a \triangleleft_{t+s}b$$
hence $F(A)$ satisfies the first property. Next, we see that:

$$ (a \triangleleft_t b ) \triangleleft_s c =e^{s [\cdot,c]}e^{t [\cdot,b]}a=e^{s [\cdot,c]}e^{t [\cdot,b]}e^{-s [\cdot,c]}e^{s [\cdot,c]}a $$ $$=e^{t e^{s [\cdot,c]}[\cdot,b]e^{-s [\cdot,c]}} (a\triangleleft_s c) =e^{t [\cdot,(b\triangleleft_s c)]} (a\triangleleft_s c) =(a\triangleleft_s c) \triangleleft_t (b\triangleleft_s c) $$

Note that if $A$ was a Lie algebra, so that $[a,a]=0$ for all $a \in A$, then we would also have that:
$$ a \triangleleft_s a = e^{s[\cdot,a]}a = a  $$
making $F(A)$ a Lie quandle in this case.


We note that for $F(A)=F(B) =(X,\triangleleft)$, we must have that $A$ and $B$ share the same underlying set of vectors $X$. We also find that for all $x,y \in X$:

$$ [x,y]_A = \frac{d}{ds}\bigg|_{s=0} (x \triangleleft_s y) =   [x,y]_B$$
Hence $A=B$, and so $A \mapsto F(A)$ is an injective function from the set of finite dimensional real Leibniz algebras to the set of smooth $\mathbb{R}$-families of smooth racks. 

 For a Leibniz algebra homomorphism $f:A \to B$, we have that:
$$ f(a\triangleleft^A_s b) = f\big(  e^{s  [\cdot,b]_A}(a)\big)= f\left(  \sum_{j=0}^{\infty} \frac{s^j}{j!} ([\cdot,b]_A)^j(a)\right)$$
$$=  \sum_{j=0}^{\infty} \frac{s^j}{j!} f\left(  ([\cdot,b]_A)^j(a)\right)=  \sum_{j=0}^{\infty} \frac{s^j}{j!}   ([\cdot,f(b)]_B)^j(f(a)) $$ $$=  e^{s [\cdot,f(b)]_B}(f(a)) = f(a)\triangleleft^B_sf(b) $$
Hence $F(f)=f$ (which clearly preserves identities and compositions) leads to $F$ being a functor from the category of finite dimensional real Leibniz algebras to the category of smooth $\mathbb{R}$-families of smooth racks. Moreover, $F$ is manifestly faithful, and so $F$ is a functor that is injective on both objects and faithful.
\end{proof}

\vspace{5mm}

This categorical embedding is telling us that the category of Leibniz racks has the same relationship to the category of finite-dimensional real Leibniz algebras that the category of finite dimensional $\mathbb{R}$-vector spaces and $\mathbb{R}$-linear maps has to the category of smooth manifolds and smooth maps. The same is true of Lie quandles with respect to finite-dimensional real Lie algebras, as Fritz \cite{Fritz} suggested. 

Leibniz racks are not only interesting for their connection to Leibniz algebras, as they are in some sense the fundamental building block of all smooth $G$-families of smooth racks. We will conclude this section with evidence for the previous claim with a demonstration of the equivalence of categories between certain smooth $G$-families of smooth racks and smooth $\mathfrak{g}$-families of Leibniz racks for finite-dimensional real Lie algebras $\mathfrak{g}$, following from the Lie-Cartan theorem. 

\vspace{5mm}
We have the following proposition.

\begin{proposition}
The category of smooth (simply-connected) Lie group families of smooth racks is isomorphic to a full subcategory of the category of smooth (real, finite dimensional) Lie algebra families of Leibniz racks.
\end{proposition}

\begin{comment}
    The last portion of this argument will allow us to remove "a full subcategory of" from the claim above.
\end{comment}

\begin{proof}
 For any Lie group $G$, and any smooth $G$-family of smooth racks $(R,\triangleleft)$, we can define $F(R,\triangleleft) = (R,\triangleleft\!\!\triangleleft)$ where $\triangleleft\!\!\triangleleft: R^2\times \mathfrak{g}\times\mathbb{R} \to R$ we define:
$$ r {\triangleleft\!\!\triangleleft}_t^v r' =r \triangleleft_{\text{exp}(tv)} r' \text{ for all }r,r' \in R$$
 where $\text{exp}:\mathfrak{g} \to G$ is the exponential map. We note that:
 $$ (r_1 {\triangleleft\!\!\triangleleft}_t^u r_2 ){\triangleleft\!\!\triangleleft}_s^v r_3  = (r_1 \triangleleft_{\text{exp}(tu)} r_2) \triangleleft_{\text{exp}(sv)} r_3$$
$$= (r_1\triangleleft_{\text{exp}(sv)} r_3) \triangleleft_{\text{exp}(-sv)\text{exp}(tu)\text{exp}(sv)} (r_2\triangleleft_{\text{exp}(sv)} r_3) $$
$$= (r_1\triangleleft_{\text{exp}(sv)} r_3) \triangleleft_{\text{exp}(-sv)\text{exp}(tu)\text{exp}(sv)} (r_2\triangleleft_{\text{exp}(sv)} r_3) $$
So we note that $ t \mapsto  \text{exp}(-sv)\text{exp}(tu)\text{exp}(sv)$ has the property that:
$$ t+t' \mapsto \text{exp}(-sv)\text{exp}(tu)\text{exp}(sv)\text{exp}(-sv)\text{exp}(t'u)\text{exp}(sv) $$
$$ = \text{exp}(-sv)\text{exp}((t+t')u)\text{exp}(sv) $$
and $t=0$ implies that $\text{exp}(-sv)\text{exp}(tu)\text{exp}(sv)  =e \in G$. Hence $$t \mapsto  \text{exp}(-sv)\text{exp}(tu)\text{exp}(sv)$$ is a Lie homomorphism from $\mathbb{R}$ into $G$, and therefore it is of the form $t \mapsto \text{exp}(tw)$ for some $w \in \mathfrak{g}$. Now we note that:
 $$ d_0 (\text{exp}(-sv)\text{exp}(tu)\text{exp}(sv))(1) = d_0 (\Psi_{\text{exp}(-sv)} \circ \text{exp}(tu))(1) $$
$$= d_e (\Psi_{\text{exp}(-sv)}) \circ d_0(\text{exp}(tu))(1) = \text{Ad}_{\text{exp}(-sv)}(u) = e^{-s\; \text{ad}_v}(u) $$
and hence $\text{exp}(-sv)\text{exp}(tu)\text{exp}(sv) = \text{exp}(t e^{-s \; \text{ad}_v}(u) ) $, therefore we find that:
$$ (q_1 {\triangleleft\!\!\triangleleft}_t^u q_2 ){\triangleleft\!\!\triangleleft}_s^v q_3 =  (q_1\triangleleft_{\text{exp}(sv)} q_3) \triangleleft_{\text{exp}(t e^{-s \; \text{ad}_v}(u) ) } (q_2\triangleleft_{\text{exp}(sv)} q_3)  $$
$$=  (q_1{\triangleleft\!\!\triangleleft}^v_s q_3) \triangleleft_{t   }^{e^{-s \; \text{ad}_v}(u)} (q_2 {\triangleleft\!\!\triangleleft}^v_s q_3 )  $$

Now for any particular $v \in \mathfrak{g}$, the above tells us that:

$$ (q_1 {\triangleleft\!\!\triangleleft}_t^v q_2 ){\triangleleft\!\!\triangleleft}_s^v q_3 = (q_1{\triangleleft\!\!\triangleleft}^v_s q_3) \triangleleft_{t   }^{e^{-s \; \text{ad}_v}(v)} (q_2 {\triangleleft\!\!\triangleleft}^v_s q_3 )  $$ $$= (q_1{\triangleleft\!\!\triangleleft}^v_s q_3) \triangleleft_{t   }^{ v} (q_2 {\triangleleft\!\!\triangleleft}^v_s q_3 )  $$
while we can also easily calculate that:
$$  q  {\triangleleft\!\!\triangleleft}_t^v q = q \triangleleft_{\text{exp}(tv)}  q  = q$$
and:
$$ (q_1 {\triangleleft\!\!\triangleleft}_t^v q_2){\triangleleft\!\!\triangleleft}_s^v q_2 =(q_1 \triangleleft_{\text{exp}(tv)}  q_2)\triangleleft_{\text{exp}(sv)}  q_2 $$
$$=(q_1 \triangleleft_{\text{exp}(tv)}  q_2)\triangleleft_{\text{exp}(sv)}  q_2 = q_1 \triangleleft_{\text{exp}(tv)\text{exp}(sv)} q_2   $$
$$ =  q_1 \triangleleft_{\text{exp}((t+s)v)} q_2 = q_1{\triangleleft\!\!\triangleleft}_{t+s}^v q_2 $$
Finally, we note that for all $r,r' \in R$, all $a \in A$, and all $t \in \mathbb{R}$:
$$ r \triangleleft \!\! \triangleleft_t^a r' = r \triangleleft_{\exp(ta)}r' = r \triangleleft \!\! \triangleleft_1^{ta} r'   $$
Hence, for any smooth $G$-family of smooth racks, $F(Q,\triangleleft)$ is a smooth $\mathfrak{g}$-family of Leibniz racks. Now suppose $A:Q \to Q'$ is a homomorphism from the smooth $G$-family of smooth racks $(Q,\triangleleft)$ to the smooth $G'$-family of smooth racks $(Q',\triangleleft')$, meaning that:
$$ A(q_l\triangleleft_g q_r) = A(q_l) \triangleleft_{a(g)} A(q_r) $$
for a Lie endomorphism $a$ of $G$. Define $F(A):F(Q,\triangleleft) \to F(Q',\triangleleft')$ have the same underlying function as $A:Q \to Q$ (Allowing for notational abuse, this can be stated concisely as $F(A) =A$). Now note that:
$$ F(A)(q_l {\triangleleft\!\!\triangleleft}_t^v q_r) = A(q_l \triangleleft_{\text{exp}_G(tv)} q_r ) = A(q_l) \triangleleft_{a(\text{exp}_G(tv))} A(q_r) $$ $$= A(q_l) \triangleleft_{\text{exp}_{G'}(td_ea(v))} A(q_r) =  (F(A)(q_l)) {\triangleleft\!\!\triangleleft}_{t}^{d_ea(v)} (F(A)(q_r)). $$
We note that $F(A)=A$ automatically implies that $F$'s morphism part preserves composition and identities; hence $F$ is a functor from the category of smooth Lie group families of smooth racks to the category of smooth Lie algebra families of Leibniz racks.

We note that the morphism part of $F$ being defined $F(A) =A$ implies that $F$ is a faithful functor, hence the category of smooth Lie group families of smooth racks can be considered a category which is concrete over the category of smooth Lie algebra families of Leibniz racks. 





 Let $(Q,\triangleleft)$ and $(Q',\triangleleft')$ be smooth simply-connected Lie group families of smooth racks, and let  $\phi : Q \to Q'$ be a homomorphism of smooth Lie algebra families of Leibniz racks from $F(Q,\triangleleft) $ to $F(Q',\triangleleft')$. Hence there is a Lie algebra homomorphism $\varphi: \mathfrak{g} \to \mathfrak{g}'$ such that:
 $$ \phi(q_1{\triangleleft\!\!\triangleleft}_t^v q_2 ) =  \phi(q_1){\triangleleft\!\!\triangleleft'}_t^{\varphi(v)} \phi( q_2)$$
 To each Lie algebra, we can associate a unique simply connected Lie group, we will call $\mathfrak{g}$'s Lie group $G$ (hence $(Q,\triangleleft)$ is a smooth $G$-family of smooth racks), and we will use $G'$ to refer to the Lie group associated to $\mathfrak{g}'$ (similarly, $(Q',\triangleleft')$ is a smooth $G'$-family of smooth racks). Moreover, to any Lie algebra homomorphism $\varphi:\mathfrak{g} \to \mathfrak{g}' $ there is a unique Lie group homomorphism $\Phi:G \to G'$ such that $T_e \Phi = \varphi$. For any $q_1,q_2 \in Q$ and any $g \in G$, there is a $v \in \mathfrak{g}$ such that $\text{exp}_G(v) =g$, hence:
 $$ \phi(q_2\triangleleft_g q_2) =  \phi(q_2{\triangleleft\!\!\triangleleft}_1^v q_2) =  \phi(q_2){\triangleleft\!\!\triangleleft'}_1^{\varphi(v)} \phi(q_2) $$ $$=  \phi(q_2)\triangleleft'_{\text{exp}_{G'}(\varphi(v))}\phi(q_2) = \phi(q_2)\triangleleft'_{\Phi(g)}\phi(q_2) $$
 Therefore $\phi$ is a homomorphism from $(Q,\triangleleft)$ to $(Q',\triangleleft')$ in the category of smooth Lie group families of smooth racks. Hence we have that: $$F(\phi:(Q,\triangleleft)\to (Q',\triangleleft')) = \phi:F(Q,\triangleleft)\to F(Q',\triangleleft')$$ it follows that $F$ is a full functor when restricted to the subcategory of all smooth simply connected Lie group families of smooth racks.
 
 
Consequently, if $(Q,\triangleleft)$ and $(Q',\triangleleft')$ are both smooth simply connected Lie group families of smooth racks such that $F(Q,\triangleleft) = F(Q',\triangleleft')$ and so that the identity map relating $F(Q,\triangleleft) $ and $ F(Q',\triangleleft')$ both lie in the image of $F$'s morphism part when restricted to $\text{Mor}((Q,\triangleleft),(Q',\triangleleft'))$ and $\text{Mor}((Q',\triangleleft'),(Q,\triangleleft))$, then $Q=Q'$ and $\triangleleft\!\!\triangleleft = {\triangleleft\!\!\triangleleft }'$. As we have seen from the fullness argument, this implies $(Q,\triangleleft)$ and $(Q',\triangleleft')$'s operations are parameterized by the same simply connected Lie group $G$, and furthermore that the identity map $\text{Id}:Q \to Q'$ determines a morphism from $(Q,\triangleleft)$ to $(Q',\triangleleft')$. Hence $F$ is not only fully-faithful on the subcategory of smooth simply connected Lie group families of smooth racks, it is also strictly injective on objects, and hence it determines an isomorphism of categories onto its image.

\begin{comment}
    Finally, to see that $F$'s image is the whole category of smooth (real, finite dimensional) Lie algebra families of Leibniz racks, let $(R,\triangleleft)$ be any object in this category, and let $\mathfrak{g}$ be the Lie algebra parameterizing its operations. By the Lie-Cartan theorem, there is a unique simply connected (real) Lie group $G$, and every $g \in G$ can be written as $\exp(u_g)\exp(v_g) =g$ for some $u_g,v_g \in \mathfrak{g}$. Define the operation $\triangleleft': R^2 \times G \to R$ via:

$$ p \triangleleft'_{g}q = (p \triangleleft_{(u_g,1)} q )\triangleleft_{(v_g,1)} q $$



(Show that this is well-defined)



Note that $e = \exp(0)\exp(0)$ implies that for each $r,r' \in R$:
$$ (p \triangleleft_e' q) = (p \triangleleft_{(0,1)} q )\triangleleft_{(0,1)} q =(p \triangleleft_{(0,0)} q )\triangleleft_{(0,0)} q = p   $$
On the other hand, for any $p \in R$ and any $g \in G$ we always have:
$$ (p \triangleleft_g p) = (p \triangleleft_{(u_g,1)} p )\triangleleft_{(u_h,1)} p  = p \triangleleft_{(u_g,1)} p = p$$
Finally, for any $r,r',r'' \in R$ and any $g,h \in G$ we have $$gh = \exp(u_g)\exp(v_g)\exp(u_h)\exp(v_h),$$ and hence:

$$( p \triangleleft'_g q) \triangleleft'_h r = (((p \triangleleft_{(u_g,1)} q) \triangleleft_{(v_g,1)} q) \triangleleft_{(u_h,1)} r) \triangleleft_{(v_h,1)} r $$


(Finish generalized right distributivity, group law is a special case when $q=r$)

Now we can conclude that $(R,\triangleleft')$ is a smooth $G$-family of smooth racks, and now by considering the operation $\triangleleft\!\!\triangleleft'$ that $\triangleleft'$ is sent to under $F$, we see that:

$$ r \triangleleft\!\!\!\triangleleft'_{(v,t)} r' = r \triangleleft'_{\exp(tv)} r'  = (r \triangleleft_{(tv,1)} r') \triangleleft_{(0,1)} r' =r \triangleleft_{(tv,1)} r' = r \triangleleft_{(v,t)} r' $$

Hence $F$'s restriction to the subcategory whose operations are parameterized by simply connected Lie groups is surjective on objects, and hence combined with the properties we've previously demonstrated we see that $F$'s restriction to this category is an isomorphism of categories.

\end{comment}


\end{proof}
 
\vspace{5mm}

 
 Hence we can see explicitly that any smooth $G$-family of smooth racks (smooth $G$-family of smooth quandles) can be reduced to a collection of Leibniz racks (Lie quandles), parallel to the manner in which simply connected Lie groups can be reconstituted from their Lie algebras.




 \section{Classification of smooth $\mathbb{R}^n$-families of Alexander quandles on $\mathbb{R}^m$} \label{Sec4}


 Consider the smooth $\mathbb{R}^n$-family of Alexander quandles on $\mathbb{R}^m $, that being the pairs $(\mathbb{R}^m,\triangleleft)$ where the operation $\triangleleft$ is given by:
$$ x \triangleleft_y x' = \alpha_y(x) + (\text{Id} - \alpha_y)(x') $$
for a function $y \mapsto \alpha_y \in GL_m(\mathbb{R}) $ for each $y \in \mathbb{R}^n$. Then from the group law axiom for smooth $G$-families of smooth quandles, we find that:
$$ \alpha_{y+y'}(x) + (\text{Id} - \alpha_{y+y'})(x')$$ $$ = x \triangleleft_{y+y'} x' =  (x \triangleleft_yx' ) \triangleleft_{y'} x' $$ $$ =  \alpha_{y'}(x \triangleleft_{y} x' ) + (\text{Id} - \alpha_{y'})(x') $$
$$ =  \alpha_{y'} \alpha_y(x) + (\alpha_{y'} - \alpha_{y'}\alpha_y)(x')   + (\text{Id} - \alpha_{y'})(x') $$
$$ =  \alpha_{y'} \alpha_y(x) + (\text{Id}  - \alpha_{y'}\alpha_y)(x')    $$
$$ \implies \alpha_{y'} \alpha_{y} = \alpha_{y+y'} $$
Hence $\alpha_y =  e^{ \sum_{j=1}^n y_j v_j }$ for some collection $v=\{v_j\}_{j=1}^n \subset \mathfrak{gl}_m(\mathbb{R})$ of pairwise commuting $m$ by $m$ matrices:
$$ x \triangleleft_y^v x' =  e^{   \sum_{j=1}^n y_j v_j}(x) + (\text{Id} -  e^{  \sum_{j=1}^n y_j v_j })(x') $$
Now suppose that there was an isomorphism $\varphi:\mathbb{R}^m\to \mathbb{R}^m$ from the Lie quandle $(\mathbb{R}^m,\triangleleft^v)$ to the Lie quandle $(\mathbb{R}^m,\triangleleft^{u})$. Then $\varphi$ would be a diffeomorphism such that:

%%%%%%%%%%%%%%%%%%%%%%
$$  e^{ \sum_{j=1}^ny_j u_j  }(\varphi(x)) + (\text{Id} - e^{\sum_{j=1}^ny_j u_j })(\varphi(x')) $$ $$= \varphi(x)  \triangleleft_y^u \varphi(x') = \varphi(x \triangleleft_y^v x')  = \varphi(e^{\sum_{j=1}^ny_j v_j }(x) + (\text{Id} - e^{\sum_{j=1}^ny_j v_j })(x') ) $$
$$  \implies  e^{ \sum_{j=1}^ny_j u_j }(\varphi(x)) + (\text{Id} - e^{\sum_{j=1}^ny_j u_j })(\varphi(x')) $$ $$= \varphi(e^{\sum_{j=1}^ny_j v_j }(x) + (\text{Id} - e^{\sum_{j=1}^ny_j v_j })(x') ) $$
%%%%%%%%%%%%%%%%%%%%%%%%%%%%

$$  d_0(e^{ \sum_{j=1}^ny_j u_j }(\varphi(x)) + (\text{Id} - e^{\sum_{j=1}^ny_j u_j })(\varphi(x')))( \partial_{y_j} ) =   u_j (\varphi(x) -  \varphi(x')) $$
$$ d_0(\varphi(e^{\sum_{j=1}^ny_j v_j }(x) + (\text{Id} - e^{\sum_{j=1}^ny_j v_j })(x') ) )(\partial_{y_j})  =(d_x\varphi \circ v_j) (x   -  x')  $$
$$ \partial_{y_j}  \implies  u_j(\varphi(x) -  \varphi(x')) =  (d_x\varphi\circ  v_j)(x   -  x')  $$
$$    u_j(\varphi(x) -   \varphi(x) - d_x\varphi(x'-x) - o(x'-x)) =  d_x\varphi\circ  v_j(x   -  x')  $$
 $$  \implies   (u_j\circ  d_x\varphi )(x-x') + o(x-x') = ( d_x\varphi\circ  v_j)(x   -  x')  $$
$$  \implies    o(x-x') =    (d_x\varphi\circ v_j -u_j \circ  d_x\varphi)(x-x') $$
$$  \implies   (u_j \circ  d_x\varphi )(x-x') =    (d_x\varphi\circ v_j)(x-x') $$
Because $x'$ is arbitrary, we must have that $u_j \circ  d_x\varphi = d_x\varphi\circ v_j$ for each index $j$, and hence $u_j$ and $v_j$ are conjugate to each other by the same invertible linear map for all $j=1,...,n$.

 
 On the other hand, let $A$ be an invertible linear transformation such that $u_j = Av_jA^{-1}$ for pairwise commuting sets $\{v_j\}_{j=1}^n,\{u_j\}_{j=1}^n$ of $m$ by $m$ matrices, and define $\varphi(x) = Ax$. Then $\varphi$ is a diffeomorphism with:
 $$ \varphi (x\triangleleft_y^vx' ) = A\circ e^{\sum_{j=1}^ny_j v_j}(x)+(A-A\circ e^{\sum_{j=1}^ny_j v_j})(x') $$
 $$= A\circ e^{\sum_{j=1}^ny_j v_j}((A^{-1}\circ A )(x))+(A -A \circ e^{\sum_{j=1}^ny_j v_j})( (A^{-1} \circ A )(x')) $$
  $$= A \circ e^{\sum_{j=1}^ny_j v_j}\circ A^{-1}  (\varphi(x))+(\text{Id}-A \circ e^{\sum_{j=1}^ny_j v_j}\circ A^{-1})(  \varphi(x')) $$
  $$=  e^{A  \sum_{j=1}^ny_j v_j A^{-1}}  (\varphi(x))+(\text{Id}- e^{  A \sum_{j=1}^ny_j v_j A^{-1}})(  \varphi(x'))  $$
   $$=  e^{\sum_{j=1}^ny_j u_j}  (\varphi(x))+(\text{Id}- e^{ \sum_{j=1}^ny_j u_j})(  \varphi(x')) = \varphi(x) \triangleleft_y^u \varphi(x') $$
   Hence we have that $(\mathbb{R}^m,\triangleleft^v)$ is isomorphic to $(\mathbb{R}^m,\triangleleft^u)$ as a smooth $\mathbb{R}^n$-family of smooth quandles if and only if the vectors $v_j$ and $u_j$ are related by conjugation by the same invertible matrix for each $j =1,...,n$.
   
    \section{Noether's First Theorem} \label{Sec5}



Fritz proposed a generalization of Noether's first theorem \cite{Fritz} , and that potentially connectedness may play some role as a hypothesis for this general version of the theorem on Lie quandles. We'll begin this section by demonstrating that connectedness is not a necessary condition to imply that a Lie quandle obeys an analogue of Noether's first theorem, though we say nothing about suffficiency here. 

We will begin by recalling the form of Noether's first theorem in the context of Lie quandles, which is that a Lie quandle $(Q,\triangleleft)$ satisfies:
$$    q \triangleleft_t p = q \text{ for all } t \in \mathbb{R}    \iff  p \triangleleft_t q = p \text{ for all } t \in \mathbb{R} ,  \text{ for all }q,p \in Q  $$
This theorem is true of the natural Lie quandle associated to a Lie algebra, which motivated Fritz to ask what properties a general Lie quandle must have to obey it. This theorem has an obvious generalization to the situation of a smooth $G$-family of smooth racks $(R,\triangleleft)$:
 $$  q \triangleleft_g r = q \text{ for all } g \in G   \iff  r \triangleleft_g q = r \text{ for all } g \in G ,  \text{ for all }r,q \in R   $$
It is clear already that we will need additional hypotheses on a smooth $G$-family of smooth racks in order for it to obey Noether's first theorem, since not all Leibniz algebras possess this property. The necessary and sufficient conditions for a Leibniz algebra to obey Noether's first theorem are given in the proposition below.


\vspace{5mm}


\begin{proposition}
The smooth $\mathbb{R}$-family of smooth racks associated to a (finite dimensional, real) Leibniz algebra $A$ obeys Noether's first theorem if and only if $[a,b]=0$ implies that $[b,a]=0$ for all $a,b \in A$.    
\end{proposition}

\begin{proof}
     First, suppose that $[a,b]=0$ implies that $[b,a]=0$ for all $a,b \in A$, and that $e^{t[\cdot,b]}(a) = a \triangleleft_t b = a $ for all $ t \in \mathbb{R} $. Then we must have that:
$$ 0 = \frac{d}{dt}\bigg|_{t=0}e^{t[\cdot,b]}(a)  = [a,b]  $$
hence $[b,a]=0$ by hypothesis. This then implies that:
$$ b \triangleleft_t a = e^{t[\cdot,a]}(b) = \sum_{j=0}^\infty \frac{t^j}{j!} [\cdot,a]^j(b) = b  $$
For the other direction, suppose $A$ obeys Noether's first theorem and $[a,b] =0$. Then it follows that $ a \triangleleft_t b = a $ for all $t \in \mathbb{R}$, and so $b\triangleleft_t a =b$ for all $t \in \mathbb{R}$. Thus we have that:
$$ 0 = \frac{d}{dt}\bigg|_{t=0}e^{t[\cdot,a]}(b) = [b,a]. $$\end{proof}

\vspace{5mm}


The property above, that $[a,b]=0$ implies that $[b,a]=0$ for all $a,b \in A$, does not force the bracket of our Leibniz algebra to be either symmetric or anti-symmetric, which can be seen from the existence of nontrivial Lie brackets, and from the Leibniz algebra structure on $\mathbb{R}^2$ given by bilinearly extending $[e_1,e_1]=e_2$, $[e_1,e_2]=[e_2,e_1]=[e_2,e_2]=0$. Another way in which Fritz conjecture can be safely modified is by looking for a hypothesis weaker than connectedness, as evidenced by the following proposition.





\vspace{5mm}

   
 
  
   \begin{theorem}
       
   All smooth $G$-families of smooth quandles of the form:
   $$ f'H\triangleleft_{g} fH = f\rho_{g}(f^{-1}f')H  $$
   where $F$ is a Lie group, $H \subset F$ is a Lie subgroup, and $g \mapsto \rho_g \in \text{Aut}(F)$ is a $G$-action on $F$ such that $\rho_g(H) \subset H$ and $f^{-1}\rho_{g} (f) \in H \implies f \rho_{g}(f^{-1}) \in H$ for each $g \in G$ have the first Noether property (In particular, for every Lie group $F$ with $H=\{e\}$ and every Lie group action $ G \ni g \mapsto \rho_g:F \to F $).  
   \end{theorem}

   
\begin{proof}
    Suppose that $f'H\triangleleft_{g} fH = f'H$, meaning that:
$$ f\rho_{g}(f^{-1}f') H = f'H   \implies  ((f')^{-1}f)^{-1}\rho_{g}( f^{-1}f' )  \in  H  $$
Now by hypothesis, we must have that:
$$   ( f^{-1}f')^{-1}\rho_{f}( (f')^{-1}f )  \in  H   \implies  f'\rho_{g}((f')^{-1}f)H = fH  $$
$$ \implies  fH\triangleleft_{g} f'H = fH$$.
\end{proof} 

\vspace{5mm}
   
Due to the existence of disconnected Lie groups such as the multiplicative group $\mathbb{R}\backslash\{0\}$, the theorem above shows that connectedness is at best a sufficient condition to imply Noether's first theorem. We also would now like to demonstrate the sufficiency of a different condition known in the quandle/rack literature as being "faithful".

\vspace{5mm}


\begin{definition}
    
A rack $(R,\triangleleft)$ is said to be \textbf{faithful} if the map defined $ r \mapsto R_r $ (where $R_r(r') = r' \triangleleft r$) for all $r \in R$ is injective.
\end{definition}

\vspace{5mm}


\begin{proposition}
Suppose $(R,\triangleleft)$ is a smooth $G$-family of smooth racks having the property that $(R,\triangleleft_{g'})  $ is faithful for some $g' $ belonging to $G$'s center, then $(R,\triangleleft)$ has the first Noether property.
\end{proposition}

\begin{proof}
     First, for a bit of notation, we write $ R^g_y(x) = x \triangleleft_g y $. Generalized right distributivity and $g'$ in the center of $G$ buys us that:
$$ (a \triangleleft_g x)\triangleleft_{g'} y = (a \triangleleft_{g'} y)\triangleleft_{g} (x \triangleleft_{g'} y) $$
$$ \implies  R^{g'}_y \circ R^g_x = R^{g}_{x \triangleleft_{g'} y}  \circ R^{g'}_y $$
While if $ x \triangleleft_{g}y =x$ for all $g \in G$ then we must also have that $ R_{x}^{h} = R_{x \triangleleft_{g}y}^{h} $  for all $g,h \in G  $. Hence we now see that:

$$ (a \triangleleft_{g'} y ) \triangleleft_{g} x = (a \triangleleft_{ g } x) \triangleleft_{g'} (y \triangleleft_{g} x) $$
$$ \implies   R_{x}^{g} \circ R^{g'}_y = R^{g'}_{y \triangleleft_g x} \circ R_{x}^g $$
While our previous observations lead to:
$$ R_{x}^{g} \circ R^{g'}_y  = R_{x\triangleleft_{g'} y}^{g} \circ R^{g'}_y =R^{g'}_y \circ R^g_x  $$
$$ R^g_x \text{ invertible} \implies R^{g'}_y  = R^{g'}_{y \triangleleft_g x} $$ $$  (R,\triangleleft_{g'}) \text{ faithful} \implies y \triangleleft_g x  =y \text{ for all }g \in G$$

\end{proof}
 


\vspace{5mm}


It is worth noting that faithfulness for a central rack operation is also not a necessary condition, since every trivial smooth $G$-family of smooth racks obeys Noether's first theorem but any having more than one point (such as the trivial structure on $\mathbb{R}$) is not faithful for any group element.

%We now move on to a classification result, concerning one particularly simple class of examples: The $\mathbb{R}^n$ families of Alexander-Lie quandles on $\mathbb{R}^m$.

\vspace{5mm}

 
\section{Discussion}


We've seen that none of the available hypotheses for Noether's first theorem can be necessary, and hence more work is needed to find sharp hypotheses. Both this article and Fritz's article leave the topic of Noether's second theorem, and its generalization to this context, unaddressed. Thorough investigation of Noether's second theorem in the context of Lie algebras and Lie quandles will be the subject of future work.

This article focuses mostly on the connection smooth $G$-families of smooth racks/quandles have to Fritz's work. Consequently, it has left further exploration of the extension of many knot invariants associated to $G$-families of quandles to the smooth setting to future work.

Finally, the small classification case dealt with here can be greatly strengthened. To that end, a sequel to this article will work towards a systematic classification of smooth $G$-families of smooth racks.

\section*{Acknowledgement} 
	The authors would like to thank Masahico Saito for his insight regarding $G$-families of quandles, as well as Tobias Fritz and Luc Ta for their helpful comments and questions.
    
%\section{References}

\begin{thebibliography}{99}
		%\setlength{\itemsep}{-2pt}

       \bib{And-Grana}{article}{
		author={Andruskiewitsch, Nicol\'{a}s},
		author={Gra\~{n}a, Mat\'{i}as},
		title={From racks to pointed Hopf algebras},
		journal={Adv. Math.},
		volume={178},
		date={2003},
		number={2},
		pages={177--243},
		
	}

 

         \bib{BES}{book}{
   author={Bardakov, Valeriy},
   author={Elhamdadi, Mohamed},
   author={Singh, Mahender},
   title={Algebra of the Yang-Baxter Equation--Skew Braces, Quandles and Cohomology},
   series={arXiv:2506.23175},
   date={2025, to appear in Springer Monographs in Mathematics},
   pages={x+338},
}

\bib{Brieskorn}{article}{
   author={Brieskorn, E.},
   title={Automorphic sets and braids and singularities},
   conference={
      title={Braids},
      address={Santa Cruz, CA},
      date={1986},
   },
   book={
      series={Contemp. Math.},
      volume={78},
      publisher={Amer. Math. Soc., Providence, RI},
   },
   isbn={0-8218-5088-1},
   date={1988},
   pages={45--115},
   
}


\bib{CCES1}{article}{
			author={Carter, J. Scott},
			author={Crans, Alissa S.},
			author={Elhamdadi, Mohamed},
			author={Saito, Masahico},
			title={Cohomology of categorical self-distributivity},
			journal={J. Homotopy Relat. Struct.},
			volume={3},
			date={2008},
			number={1},
			pages={13--63},
			
		}
		
		\bib{CCES2}{article}{
			author={Carter, J. Scott},
			author={Crans, Alissa S.},
			author={Elhamdadi, Mohamed},
			author={Saito, Masahico},
			title={Cohomology of the adjoint of Hopf algebras},
			journal={J. Gen. Lie Theory Appl.},
			volume={2},
			date={2008},
			number={1},
			pages={19--34},
			
		}

        \bib{CCEKS}{article}{
			author={Carter, J. Scott},
			author={Crans, Alissa S.},
			author={Elhamdadi, Mohamed},
			author={Karadayi, Enver},
			author={Saito, Masahico},
			title={Cohomology of Frobenius algebras and the Yang-Baxter equation},
			journal={Commun. Contemp. Math.},
			volume={10},
			date={2008},
			number={suppl. 1},
			pages={791--814},
			
		}

\bib{Covez}{article}{
   author={Covez, Simon},
   title={The local integration of Leibniz algebras},
   language={English, with English and French summaries},
   journal={Ann. Inst. Fourier (Grenoble)},
   volume={63},
   date={2013},
   number={1},
   pages={1--35},
  
}

\bib{EFT}{article}{
   author={Elhamdadi, Mohamed},
   author={Fernando, Neranga},
   author={Tsvelikhovskiy, Boris},
   title={Ring theoretic aspects of quandles},
   journal={J. Algebra},
   volume={526},
   date={2019},
   pages={166--187},
   
}

\bib{Elhamdadi}{article}{
   author={Elhamdadi, Mohamed},
   title={Distributivity in quandles and quasigroups},
   conference={
      title={Algebra, geometry and mathematical physics},
   },
   book={
      series={Springer Proc. Math. Stat.},
      volume={85},
      publisher={Springer, Heidelberg},
   },
   date={2014},
   pages={325--340},
   %review={\MR{3275946}},
   %doi={10.1007/978-3-642-55361-5_19},
}

\bib{EM}{article}{
   author={Elhamdadi, Mohamed},
   author={Moutuou, El-Ka\"{\i}oum M.},
   title={Foundations of topological racks and quandles},
   journal={J. Knot Theory Ramifications},
   volume={25},
   date={2016},
   number={3},
   pages={1640002, 17},
}


    \bib{EN}{book}{
   author={Elhamdadi, Mohamed},
   author={Nelson, Sam},
   title={Quandles---an introduction to the algebra of knots},
   series={Student Mathematical Library},
   volume={74},
   publisher={American Mathematical Society, Providence, RI},
   date={2015},
    pages={x+245},

}

\bib{ESZ}{article}{
   author={Elhamdadi, Mohamed},
   author={Saito, Masahico},
   author={Zappala, Emanuele},
   title={Continuous cohomology of topological quandles},
   journal={J. Knot Theory Ramifications},
   volume={28},
   date={2019},
   number={6},
   pages={1950036, 22},
   
}


\bib{Fritz}{article}{
   author={Fritz, Tobias},
   title={Self-distributive structures in physics},
   journal={Internat. J. Theoret. Phys.},
   volume={64},
   date={2025},
   number={3},
   pages={Paper No. 73, 10},
  
}

\bib{IIJO}{article}{
   author={Ishii, Atsushi},
   author={Iwakiri, Masahide},
   author={Jang, Yeonhee},
   author={Oshiro, Kanako},
   title={A $G$-family of quandles and handlebody-knots},
   journal={Illinois J. Math.},
   volume={57},
   date={2013},
   number={3},
   pages={817--838},
   
}

\bib{Joyce}{article}{
   author={Joyce, David},
   title={A classifying invariant of knots, the knot quandle},
   journal={J. Pure Appl. Algebra},
   volume={23},
   date={1982},
   number={1},
   pages={37--65},
   
}

\bib{Kinyon}{article}{
   author={Kinyon, Michael K.},
   title={Leibniz algebras, Lie racks, and digroups},
   journal={J. Lie Theory},
   volume={17},
   date={2007},
   number={1},
   pages={99--114},
   
}

\bib{Loday}{article}{
   author={Loday, Jean-Louis},
   title={Une version non commutative des alg\`ebres de Lie: les alg\`ebres
   de Leibniz},
   language={French},
   journal={Enseign. Math. (2)},
   volume={39},
   date={1993},
   number={3-4},
   pages={269--293},
   
}


\bib{Matveev}{article}{
   author={Matveev, S. V.},
   title={Distributive groupoids in knot theory},
   language={Russian},
   journal={Mat. Sb. (N.S.)},
   volume={119(161)},
   date={1982},
   number={1},
   pages={78--88, 160},
   
}

\bib{Rubinsztein}{article}{
   author={Rubinsztein, Ryszard L.},
   title={Topological quandles and invariants of links},
   journal={J. Knot Theory Ramifications},
   volume={16},
   date={2007},
   number={6},
   pages={789--808},
   
}

\bib{LucTa}{article}{
   author={Ta, Luc},
   title={From affine algebraic racks to Leibniz algebras and Yang-Baxter operators},
   journal={	arXiv:2512.15909,}
   %volume={28},
   date={2026},
}


\bib{Takasaki}{article}{
   author={Takasaki, Mituhisa},
   title={Abstraction of symmetric transformations},
   language={Japanese},
   journal={T\^{o}hoku Math. J.},
   volume={49},
   date={1943},
   pages={145--207},
  
}


	

        \end{thebibliography}


\end{document}
